\chapter{Operational states}
\label{ch:states}

This chapter derives the operational regimes of the train and validates them. The regimes
are recovered by clustering, and, in line with the feature-role split
(Section~\ref{sec:roles}), they are defined from motion alone. Only the kinematic window
summaries (speed, acceleration, jerk, and their spreads) enter the clustering; the pneumatic
sensors are deliberately withheld so that they remain independent evidence in the chapters
that follow.

\section{Clustering into operational states}
\label{sec:kselection}

The states were obtained by k-means clustering of the motion window features, with the
number of clusters chosen by internal validation rather than fixed in advance. Two clustering
methods (k-means and hierarchical clustering on a sample) were run for $k$ from 2 to 6, and
four indicators were compared (Figure~\ref{fig:kselection}). The \emph{silhouette} score~\cite{Rousseeuw:87}
measures how well each window fits its own cluster relative to the nearest alternative; it
ranges over $[-1, 1]$, where values near 1 indicate well-separated clusters, values near 0 an
ambiguous assignment, and higher is better. The \emph{Davies--Bouldin} index~\cite{DaviesBouldin:79}
measures the average overlap between a cluster and its most similar neighbour; it is bounded
below by 0 and unbounded above, with lower values better. The \emph{Bayesian information
criterion} (BIC)~\cite{Schwarz:78} of a Gaussian mixture balances goodness of fit against model
complexity; it is not bounded to a fixed range and its absolute value is not interpretable, so
only differences between models are read, with lower values better. The \emph{adjusted Rand
index} (ARI) measures agreement between two clusterings; it takes the value 1 for identical
partitions and about 0 for the agreement expected by chance (it can be slightly negative), and
is used here both to compare the k-means and hierarchical labels and to assess stability under
resampling.

The per-$k$ values of all four indicators are collected in Table~\ref{tab:kselection}. The
silhouette score (Figure~\ref{fig:kselection}a) rises monotonically with $k$, from 0.48 at
$k=2$ to 0.60 at $k=6$, and the Davies--Bouldin index (Figure~\ref{fig:kselection}b) likewise
keeps improving as clusters are added, from 1.20 at $k=3$ to 0.95 at $k=6$. Both therefore
reward ever finer partitions and neither picks out a value on its own; this is the expected
behaviour of internal compactness indices, so the choice cannot rest on them alone. The two
clustering methods agree closely at both of the small $k$ that were cross-checked (cross-method
ARI 0.81 at $k=3$ and 0.77 at $k=4$), and the $k=4$ solution is stable under resampling
(bootstrap ARI about 0.87).

The BIC (Figure~\ref{fig:kselection}c) is informative for a specific reason: unlike the
silhouette and Davies--Bouldin scores it explicitly penalises model complexity, so rather than
improving without bound it is expected to flatten once additional clusters stop describing real
structure. Its absolute value is not meaningful; only the size of each successive drop is. The
largest single drop is the $k=3\rightarrow4$ step ($-121{,}900$), but the $k=4\rightarrow5$ step
is nearly as large ($-110{,}900$), so the BIC does not by itself separate four clusters from
five. What it does show is where the marginal gain collapses: the $k=5\rightarrow6$ step
($-54{,}000$) is less than half of either earlier step. The BIC therefore rules out large $k$
but leaves the choice between four and five to be settled on other grounds.

That choice is made on parsimony and physical interpretability. None of the internal indices
favours five clusters over four: the silhouette and Davies--Bouldin scores keep rising for any
$k$, and the BIC does not distinguish the two. The smaller model is therefore preferred unless
a fifth cluster earns its place with a distinct operating mode, and it does not. Four clusters
already exhaust the stop--accelerate--cruise--brake cycle of a metro duty cycle
(Table~\ref{tab:states}), each with its own signature in speed and the sign of acceleration,
whereas a fifth group adds no regime that is not already present. Four clusters is kept as the
most parsimonious description consistent with the metrics and with the physics of the duty
cycle.

\begin{table}[htbp]
  \centering
  \caption{Cluster-count selection metrics per $k$. Silhouette (higher better, range
  $[-1,1]$), Davies--Bouldin (lower better, $\geq 0$), cross-method ARI ($k$-means versus
  hierarchical; 1 is identical, computed for the two small $k$ that were cross-checked), and the
  Gaussian-mixture BIC (lower better; only differences are meaningful). $\Delta$BIC is the drop
  from the previous $k$. The largest drop is at $k=3\rightarrow4$, but $k=4\rightarrow5$ is
  comparable; the gain collapses only after $k=5$.}
  \label{tab:kselection}
  \begin{tabular}{rrrrrr}
    \toprule
    $k$ & Silhouette & Davies--Bouldin & Cross-method ARI & BIC ($\times 10^{3}$) & $\Delta$BIC ($\times 10^{3}$) \\
    \midrule
    2 & 0.476 & 0.999 & n/a  & $-999.1$  & n/a      \\
    3 & 0.535 & 1.196 & 0.81 & $-1088.8$ & $-89.8$  \\
    4 & 0.567 & 1.091 & 0.77 & $-1210.8$ & $-121.9$ \\
    5 & 0.587 & 0.990 & n/a  & $-1321.6$ & $-110.9$ \\
    6 & 0.605 & 0.945 & n/a  & $-1375.7$ & $-54.0$  \\
    \bottomrule
  \end{tabular}
\end{table}

\begin{figure}[htbp]
  \centering
  \includegraphics[width=0.9\linewidth]{kinematic_kselection.png}
  \caption{Cluster-count selection on the motion features. \textbf{(a)} Silhouette rises
  monotonically and \textbf{(b)} Davies--Bouldin keeps improving, so neither compactness index
  picks a $k$ on its own. \textbf{(c)} The Gaussian-mixture BIC has its largest drop at
  $k=3\rightarrow4$, but the $k=4\rightarrow5$ drop is comparable; the marginal gain collapses
  only after $k=5$. Four groups is retained on grounds of parsimony and physical
  interpretability (see text and Table~\ref{tab:kselection}).}
  \label{fig:kselection}
\end{figure}

\label{sec:fourstates}
Table~\ref{tab:states} numbers the four clusters and, for each, gives its held-out share, its
median motion, and the reasoning that maps it onto an operating regime. The regime labels used
in the rest of this report are read off from these medians rather than assigned in advance, so
that the naming follows from the data.

\begin{table}[htbp]
  \centering
  \caption{The four clusters on the held-out window, numbered in duty-cycle order. Shares and
  medians are computed on ten-second windows; velocity and acceleration are in normalised units.
  The final column states why each cluster's median motion identifies it with the named regime.}
  \label{tab:states}
  \begin{tabularx}{\linewidth}{c r r r >{\RaggedRight}X}
    \toprule
    Cluster & Median vel. & Median accel. & Share (\%) & Interpretation and justification \\
    \midrule
    1 & 0.000 & \phantom{$-$}0.000 & 48.12 & Zero speed and zero acceleration: the train is at
        rest, labelled \emph{standing}. \\[0.4ex]
    2 & 0.305 & \phantom{$-$}0.041 & 14.34 & Moderate speed with the largest positive median
        acceleration: the train is speeding up, labelled \emph{accelerating}. \\[0.4ex]
    3 & 0.842 & \phantom{$-$}0.002 & 17.27 & Highest median speed with near-zero acceleration:
        steady running, labelled \emph{cruising}. \\[0.4ex]
    4 & 0.343 & $-$0.033 & 20.26 & Moderate speed with clearly negative acceleration: the train
        is slowing down, labelled \emph{deceleration}. \\
    \bottomrule
  \end{tabularx}
\end{table}

The resulting profile matches a typical metro duty cycle. Standing dominates at about 48\% of
windows, with zero median speed and acceleration. Cruising has the highest median speed
(0.84) and near-zero acceleration. Accelerating and deceleration sit at intermediate speeds
and are separated by the sign of acceleration ($+0.041$ against $-0.033$). Deceleration
accounts for about one window in five (20.3\%), which sets the scale for the braking analysis
in Chapter~\ref{ch:braking}. Figure~\ref{fig:profiles} shows the four motion profiles, and
Figure~\ref{fig:velaccel} the velocity--acceleration cycle the train traces through them.

\begin{figure}[htbp]
  \centering
  \includegraphics[width=0.9\linewidth]{state_profiles.png}
  \caption{Motion profiles of the four states: \textbf{(a)} median velocity, \textbf{(b)} median
  acceleration, and \textbf{(c)} median jerk. The states are separated by speed and the sign and
  magnitude of acceleration; no pneumatic channel enters the definition.}
  \label{fig:profiles}
\end{figure}

\begin{figure}[htbp]
  \centering
  \includegraphics[width=0.75\linewidth]{kinematic_vel_accel_scatter.png}
  \caption{The velocity--acceleration cycle, coloured by state. The train moves anticlockwise:
  stop, accelerate, cruise, decelerate, stop. A minority of off-curve points do not fit the
  idealised cycle and motivate future work (Chapter~\ref{ch:conclusion}).}
  \label{fig:velaccel}
\end{figure}

\section{Plausibility of four states}
\label{sec:heldbrake}

An earlier version of this analysis included brake-pressure sensors in the state definition.
It systematically mislabelled a stationary train holding its brake at a station as
deceleration: the brake was applied, so the pneumatic signal resembled braking even though
the train was not moving. This is precisely the circularity that the feature-role split
(Section~\ref{sec:roles}) is designed to avoid, and it is the reason the state definition was
restricted to motion. Table~\ref{tab:heldbrake} shows how the earlier, pneumatics-contaminated
\enquote{braking} label redistributes under the motion-only definition. Of the windows the
old label called braking, about 105{,}000 are genuine deceleration, roughly 80{,}000 are in
fact acceleration, and about 47{,}000 are a stationary train holding its brake. Rather than discarding this information whether a standing train is braked is retained as a sub-label
(\texttt{standing\_braked} versus \texttt{standing\_idle}), but it no longer contaminates the
deceleration regime, which is 99.9\% genuine deceleration under the corrected definition.

\begin{table}[htbp]
  \centering
  \caption{Reassignment of the old, pneumatics-contaminated \enquote{braking} label under the
  motion-only definition. Windows counted at the window level.}
  \label{tab:heldbrake}
  \begin{tabular}{lr}
    \toprule
    Motion-only state & Windows from old \enquote{braking} \\
    \midrule
    accelerating & 80{,}226 \\
    deceleration & 104{,}908 \\
    cruising     & 519 \\
    standing     & 46{,}723 \\
    \bottomrule
  \end{tabular}
\end{table}

As a second check, the motion-only clustering was compared against a classifier trained on
the full feature set. A Random Forest predicting the state labels from both the kinematic and
the pneumatic window summaries reaches 0.994 accuracy on held-out windows, and its most
important features are the kinematic ones: the mean, maximum, and minimum of acceleration,
and the mean and spread of speed. Pneumatic channels appear only much lower in the ranking.
Even when the pneumatic signals are made available, then, the states remain separable by
motion alone, which confirms that they are motion-defined regimes rather than an artefact of
the pneumatic data.

\section{Transitions and dwell}

The states occur in the expected order (stop, accelerate, cruise, brake, stop), and the
transition matrix is dominated by this cycle (Figure~\ref{fig:transitions}). Mean dwell times
match the physical picture: standing about 54\,s (with very high variance, since layovers and
terminus stops are long and irregular), accelerating about 17\,s, cruising about 21\,s, and
deceleration about 24\,s. The state mix is stable from month to month
(Figure~\ref{fig:bymonth}), which indicates that the regimes are a property of how the train
is operated rather than of any particular period in the recording.

\begin{figure}[htbp]
  \centering
  \includegraphics[width=0.7\linewidth]{transition_diagram.png}
  \caption{State transition diagram. The dominant path is the operating cycle
  stop\,$\rightarrow$\,accelerate\,$\rightarrow$\,cruise\,$\rightarrow$\,decelerate\,$\rightarrow$\,stop.}
  \label{fig:transitions}
\end{figure}

\begin{figure}[htbp]
  \centering
  \includegraphics[width=0.9\linewidth]{state_distribution_by_month.png}
  \caption{State mix by month. The proportions are stable across the recording, so the regimes
  reflect the operating pattern rather than a seasonal artefact.}
  \label{fig:bymonth}
\end{figure}

\section{Robustness to window length}
\label{sec:window-sensitivity}

Because the ten-second window is a design choice, the clustering was re-run at 5, 10, 15, and
20 seconds as a read-only sensitivity check that left the canonical ten-second models
untouched (Figure~\ref{fig:windowsens}). The deceleration share is essentially flat across
all four sizes, about 20\% at each. The state \emph{chatter} rate, the fraction of windows
whose state differs from an adjacent window, rises with window length, from about 0.19 at
5\,s to 0.54 at 20\,s. Short windows are therefore not the source of label churn, and ten
seconds is a sensible compromise between resolution and stability.

\begin{figure}[htbp]
  \centering
  \includegraphics[width=0.9\linewidth]{window_sensitivity.png}
  \caption{Window-length sensitivity. State fractions (left) are stable across 5--20\,s; the
  deceleration share stays near 20\%. Chatter (right) rises with window length, so short
  windows are not the cause of label churn.}
  \label{fig:windowsens}
\end{figure}

\section{Pre-event state distribution}
\label{sec:preevent}

As an initial link to the maintenance record, the mix of operational states in the 24 to
72~hours before each documented failure or maintenance event was compared with the overall
baseline. Across 62 event-by-window-size chi-square tests, 61 were significant at $p<0.05$.
This significance is, however, an artefact of sample size: with 6{,}000 to 24{,}000 windows
per test, even negligible differences reach significance. The effect size is small
throughout: Cramér's $V$ does not exceed 0.07 and is typically about 0.018, an order of
magnitude below the 0.1 value conventionally regarded as small.

The shift that does exist reflects the train standing idle in the depot before scheduled work.
Ahead of planned revisions the standing fraction rises by roughly 45 to 52 percentage points.
Restricting the test to moving windows, by dropping the standing state, reduces the effect
size roughly sixfold, to a Cramér's $V$ of about 0.003, which is indistinguishable from zero.
For the eight brake failures specifically there is no such idling, since failures occur in
service, and no meaningful effect size either.

The test therefore confirms only that the pipeline registers when the train is parked before
scheduled maintenance, which is an operational sanity check rather than evidence that
pre-failure braking behaviour differs from normal operation. This is consistent with the
near-null pre-failure braking signal reported in Chapter~\ref{ch:failure}, and it motivates
examining charging and idle windows with the auxiliary sensors, rather than deceleration, in
future work.